\documentclass[conference]{IEEEtran}
\IEEEoverridecommandlockouts
% The preceding line is only needed to identify funding in the first footnote. If that is unneeded, please comment it out.
\usepackage{cite}
\usepackage{amsmath,amssymb,amsfonts}
\usepackage{algorithmic}
\usepackage{graphicx}
\usepackage{textcomp}
\usepackage{xcolor}
\def\BibTeX{{\rm B\kern-.05em{\sc i\kern-.025em b}\kern-.08em
    T\kern-.1667em\lower.7ex\hbox{E}\kern-.125emX}}
\begin{document}

\title{
\thanks{Identify applicable funding agency here. If none, delete this.}
}

\author{\IEEEauthorblockN{1\textsuperscript{st} Given Name Surname}
\IEEEauthorblockA{\textit{dept. name of organization (of Aff.)} \\
\textit{name of organization (of Aff.)}\\
City, Country \\
email address or ORCID}
\and
\IEEEauthorblockN{2\textsuperscript{nd} Given Name Surname}
\IEEEauthorblockA{\textit{dept. name of organization (of Aff.)} \\
\textit{name of organization (of Aff.)}\\
City, Country \\
email address or ORCID}}

\maketitle

\begin{abstract}
This survey paper provides a comprehensive overview of autonomous agents in Multi-Layer Large Language Models (MLLM), a rapidly evolving field that has garnered significant attention in recent years. The concept of autonomous agency in MLLM is explored through an analysis of various research papers, highlighting its potential applications, technical requirements, and theoretical foundations. The scope of this paper encompasses the introduction to large language models and agents, background and foundations of LLMs, recent advances in architecture, training, and evaluation, as well as the development, capabilities, and integration of LLM-based autonomous agents with human tasks.

Key themes and findings include the potential of LLMs in achieving human-level intelligence, their applications in social science, natural science, and engineering, and the importance of understanding user trust, agency, and performance in human-AI collaboration. The paper also discusses technical challenges and open problems in LLM and agent development, including the lack of autonomy and self-improvement in LLMs. Furthermore, it explores various applications of LLMs and agents, such as daily assistance, code development, and social simulations.

This survey contributes to the existing literature by providing a unified framework for constructing LLM-based autonomous agents, highlighting their capabilities, and discussing future directions for integrating LLMs with multimodal input, edge AI, and explainability. The paper also raises important societal implications and ethical considerations of widespread LLM and agent adoption, emphasizing the need for comprehensive regulation and evaluation. Ultimately, this survey aims to provide a thorough understanding of the state of the art in large language models and agents, shedding light on the paths forward for achieving artificial collective intelligence and autonomous agency in MLLM.
\end{abstract}

\begin{IEEEkeywords}
large language models, multimodal learning, natural language processing, machine learning, artificial intelligence
\end{IEEEkeywords}


\section{Introduction to Large Language Models and Agents}
Introduction to Large Language Models and Agents

The concept of autonomous agents in the context of large language models (LLMs) has garnered significant attention in recent years, with a plethora of research papers exploring their potential applications, technical requirements, and theoretical foundations. As highlighted in \cite{2308.11432}, LLM-based autonomous agents have diverse applications in social science, natural science, and engineering, underscoring their versatility and potential impact on various fields. A unified framework for constructing these agents has been proposed, encompassing a majority of previous work \cite{2308.11432}. Furthermore, the introduction of an open-source library called Agents \cite{2309.07870} has facilitated the development of autonomous language agents with features like planning, memory, tool usage, and multi-agent communication.

The notion of autonomy and agency in LLMs is a central theme, with discussions on whether these models can truly be considered autonomous agents \cite{2407.10735}. While some papers focus on developing autonomous agents based on LLMs \cite{2308.11432, 2411.14033}, others question the autonomy of LLMs, suggesting they lack certain necessary conditions for agency \cite{2407.10735}. The approach to achieving agency in LLMs varies, with some papers proposing the use of multi-agent systems \cite{2411.14033} and others focusing on developing more sophisticated single-agent models. For instance, \cite{2405.11106} explores LLM-based multi-agent reinforcement learning, highlighting current and future directions.

Technical details and methodologies discussed in the literature include the development of unified frameworks, open-source library implementation, and LaMAS protocols \cite{2308.11432, 2309.07870, 2411.14033}. However, limitations and challenges have also been identified, such as the unpredictability of LLMs, the gap between research and implementation, and the need for efficient resource management \cite{2407.20859, 2405.11106, 2404.13501}. These challenges underscore the importance of careful consideration and planning in the development and deployment of LLM-based autonomous agents.

The implications of LLM-based autonomous agents are far-reaching, with potential applications in various domains, including social sciences, natural sciences, and engineering \cite{2308.11432, 2411.14033}. As noted in \cite{2403.11381}, the cooperative capabilities of LLM-augmented autonomous agents can be evaluated through frameworks like Melting Pot. Moreover, the development of protocols for LaMAS involves addressing technical requirements, data privacy, and business incentives \cite{2411.14033}. The connections between LLMs, autonomy, and agency are complex and multifaceted, requiring a comprehensive understanding of the underlying theoretical foundations and technical requirements.

In conclusion, the field of LLM-based autonomous agents is rapidly evolving, with significant advancements in recent years. As research continues to uncover new developments and challenges, it is essential to maintain a nuanced understanding of the complex relationships between LLMs, autonomy, and agency. By acknowledging the key themes, limitations, and implications of this field, researchers can work towards developing more sophisticated and practical LLM-based autonomous agents that can be effectively applied in various real-world contexts \cite{2407.10735, 2411.14033}. Future research directions should focus on addressing the existing challenges, exploring new applications, and refining the theoretical foundations of LLM-based autonomous agents \cite{2308.11432, 2405.11106}.

\section{Background and Foundations: Overview of LLMs and Their Applications}
The concept of Large Language Models (LLMs) has revolutionized the field of artificial intelligence, enabling the creation of autonomous agents capable of performing complex tasks. Recent studies have explored the potential of LLMs in developing autonomous agents, leading to significant advancements in areas such as natural language processing, planning, and decision-making \cite{2308.11432}, \cite{2402.02716}. The use of LLMs as the underlying models for autonomous agents has become a recurring theme across various research papers, with many studies focusing on the development of LLM-based multi-agent systems \cite{2411.14033}, \cite{2402.01680}. These systems have shown remarkable success in complex problem-solving and world simulation, offering advantages such as dynamic task decomposition, higher flexibility, and proprietary data preservation.

Research has also led to the proposal of a unified framework for constructing LLM-based autonomous agents, encompassing a majority of previous work in the field \cite{2308.11432}. Furthermore, studies have developed a taxonomy of existing works on LLM-agent planning, categorizing them into task decomposition, plan selection, external module, reflection, and memory \cite{2402.02716}. The importance of considering practical aspects such as handling stochasticity, managing resources efficiently, and positioning research findings into categories like planning, memory, tools, and control flow when designing agentic LLMs has also been highlighted \cite{2412.04093}. Additionally, the emergence of LLM-based autonomous agents has led to a focus on the development of evaluation metrics and benchmarks to assess their performance, with various approaches being proposed \cite{2308.05960}, \cite{2409.14457}.

The ability of LLMs to plan and make decisions is a crucial aspect of developing autonomous agents, with multiple studies focusing on this area \cite{2402.02716}, \cite{2411.14033}. The development of LLM-based multi-agent systems is seen as a natural progression in the field, enabling more complex problem-solving and world simulation. However, contrasting viewpoints exist, with some studies focusing on single-LLM agent systems and others arguing that multi-LLM agent systems offer greater advantages in terms of flexibility, scalability, and task decomposition \cite{2401.03428}, \cite{2412.17481}. The surveyed papers also demonstrate a mix of technical and business perspectives, with some focusing on the technical aspects of LLM-based autonomous agents and others exploring their potential applications and benefits \cite{2411.14033}, \cite{2407.10735}.

Technical details such as LLM architecture, planning and decision-making algorithms, and evaluation metrics are critical components of developing autonomous agents \cite{2402.01680}, \cite{2308.11432}. The architecture of LLMs is a critical aspect of developing autonomous agents, with studies exploring various architectures and their implications for planning and decision-making. Researchers have proposed several algorithms and techniques for planning and decision-making in LLM-based autonomous agents, including task decomposition, plan selection, and external module integration \cite{2402.02716}, \cite{2411.14033}. However, the inherent unpredictability of LLMs poses a significant challenge in developing autonomous agents, requiring careful consideration of factors like stochasticity and resource management \cite{2412.04093}. As LLM-based autonomous agents become more complex, scalability and complexity issues arise, necessitating the development of more efficient architectures and algorithms.

The lack of standardization in the field of LLM-based autonomous agents is also a significant challenge, with different studies proposing varying approaches and methodologies, making it challenging to compare and contrast results \cite{2401.03428}, \cite{2412.17481}. Despite these limitations, the development of LLM-based autonomous agents has the potential to revolutionize various fields, including healthcare, finance, and transportation. By summarizing the key findings, common themes, contrasting viewpoints, technical details, and limitations of LLM-based autonomous agents, this analysis provides a comprehensive foundation for understanding the current state of research in this area \cite{2308.11432}, \cite{2402.02716}. The connections between LLMs, planning, and decision-making are critical to the development of autonomous agents, and further research is needed to address the challenges and limitations associated with these complex systems \cite{2411.14033}, \cite{2409.14457}. Ultimately, the development of LLM-based autonomous agents has the potential to transform the way we approach complex tasks and make decisions, leading to significant advancements in various fields.

\section{Recent Advances in LLMs: Improvements in Architecture, Training, and Evaluation}
Recent advances in Large Language Models (LLMs) have led to significant improvements in architecture, training, and evaluation, particularly in the context of autonomous agents. The potential of LLMs in achieving human-level intelligence and their application in various fields has been highlighted in several surveys \cite{2308.11432}, \cite{2402.02716}, and \cite{2411.14033}. A unified framework for constructing LLM-based autonomous agents has been proposed, encompassing a majority of previous work \cite{2308.11432}, while a taxonomy of existing works on LLM-Agent planning has been provided, categorizing them into Task Decomposition, Plan Selection, External Module, Reflection, and Memory \cite{2402.02716}.

The concept of LLM-based Multi-Agent Systems (LaMAS) has been introduced, offering advantages such as dynamic task decomposition, higher flexibility, and proprietary data preservation \cite{2411.14033}. An evaluation-driven design approach for LLM agents has been proposed, integrating online and offline evaluations to support adaptive runtime adjustments and systematic offline redevelopment \cite{2411.13768}. Furthermore, a novel paradigm of training LLM agents without modifying the LLM weights has been presented, treating functions as learnable `agent parameters' \cite{2402.11359}. These developments have significant implications for the field, enabling more efficient and effective development of autonomous agents.

The importance of evaluating and improving the performance of LLM-based autonomous agents is a common theme across the surveyed papers \cite{2308.03688}, \cite{2406.05804}, and \cite{2412.04093}. The need for flexible and adaptive architectures to support the development of LLM-based autonomous agents is emphasized, with several papers discussing the potential of LaMAS in achieving artificial collective intelligence \cite{2411.14033} and \cite{2402.01680}. The approach to training LLM agents varies, with some papers proposing modifications to the LLM weights, while others suggest treating functions as learnable `agent parameters' \cite{2402.11359}.

Technical details and methodologies employed in the surveyed papers include model training and evaluation methods \cite{2404.04442}, agent development frameworks and architectures \cite{2309.07870}, function-based training paradigms \cite{2402.11359}, and multi-agent system protocols and data privacy considerations \cite{2411.14033}. The papers also discuss the importance of human-AI collaboration and the need for fine-grained feedback from human evaluators \cite{2312.17025} and \cite{2409.02977}. However, several limitations and challenges are highlighted, including the need for more efficient and effective evaluation methods for LLM-based autonomous agents \cite{2406.00936}, the challenge of developing flexible and adaptive architectures, and the importance of addressing data privacy and security concerns in LaMAS \cite{2411.14033}.

In conclusion, recent advances in LLMs have led to significant improvements in architecture, training, and evaluation, with important implications for the development of autonomous agents. The surveyed papers highlight key themes and developments, including the potential of LaMAS, the importance of evaluation and improvement, and the need for flexible and adaptive architectures. Further research is needed to address the limitations and challenges identified, including the development of more efficient evaluation methods, flexible architectures, and secure multi-agent systems. As the field continues to evolve, it is likely that LLM-based autonomous agents will play an increasingly important role in a wide range of applications, from robotics and healthcare to finance and education \cite{2402.02716} and \cite{2411.14033}.

\section{LLM Agents: Development, Capabilities, and Integration with Human Tasks}
The development of Large Language Model (LLM) agents has been a rapidly evolving field, with significant contributions to their construction, capabilities, and integration with human tasks. A unified framework for constructing LLM-based autonomous agents has been proposed, encompassing a majority of previous work \cite{2308.11432}. This framework highlights the ability of LLM agents to perceive, control, and receive feedback from the environment to accomplish tasks autonomously, and to call external tools to ease task completion \cite{2411.14033}. The importance of integrating human tasks and knowledge into LLM agents is also emphasized, enabling them to achieve human-like decisions \cite{2402.02716}.

The analysis of existing works on LLM-agent planning reveals a taxonomy that categorizes them into Task Decomposition, Plan Selection, External Module, Reflection, and Memory \cite{2402.11359}. Furthermore, a novel paradigm for training LLM agents without modifying the LLM weights has been introduced, treating functions as learnable agent parameters \cite{2406.05804}. This approach enables more efficient and flexible training of LLM agents, and has significant implications for their development and deployment.

The use of large language models as a foundation for autonomous agents is a dominant theme in the field, with several papers highlighting their potential for achieving artificial collective intelligence \cite{2411.14033, 2402.01680}. The need for flexible and dynamic frameworks that can adapt to changing environments and tasks is also emphasized, as well as the challenges of evaluating and training LLM agents \cite{2308.03688, 2412.04093}. A unified taxonomy has been proposed to systematically review and discuss frameworks across different paradigms (tool use, planning, and feedback learning) for developing LLM-based agents \cite{2406.05804}.

The development of LLM-based multi-agent systems (LaMAS) is also a key area of research, with significant potential for achieving artificial collective intelligence \cite{2411.14033}. However, this approach also raises important questions about data privacy and security, as well as the need for more sophisticated evaluation metrics \cite{2402.01680}. The integration of human tasks and knowledge into LLM agents is critical to their development, and several papers highlight the importance of human-agent collaboration and the need for more effective interfaces \cite{2412.15701, 2308.03688}.

Despite the significant progress made in the field, there are still several limitations and challenges associated with developing LLM agents. The difficulty of achieving human-like decisions with LLM agents is a major challenge, as well as the need for more flexible and dynamic frameworks that can adapt to changing environments and tasks \cite{2402.02716, 2411.14033}. The evaluation and training of LLM agents are also key areas of research, with several papers highlighting the need for more sophisticated metrics and methods \cite{2308.03688, 2412.04093}. Overall, the development of LLM agents is a rapidly evolving field with significant potential for achieving artificial collective intelligence and improving human-agent collaboration. However, further research is needed to address the limitations and challenges associated with their development and deployment.

\section{Human-AI Collaboration: Understanding User Trust, Agency, and Performance with LLM Agents}
Human-AI collaboration has emerged as a crucial area of research, with Large Language Models (LLMs) playing a pivotal role in developing autonomous agents that can effectively collaborate with humans \cite{2308.11432}, \cite{2412.15701}. The concept of collaborative agents is built on the premise that LLMs can be leveraged to create intelligent systems that not only perform tasks autonomously but also interact and cooperate with humans seamlessly. A framework for enabling and evaluating human-agent collaboration, known as Collaborative Gym (Co-Gym), has been proposed to assess the performance of such agents \cite{2412.15701}. This framework demonstrates improved task performance when collaborative agents are used, underscoring the potential benefits of human-AI collaboration.

The versatility of LLM-based agents is a key finding across various studies, highlighting their suitability for diverse applications, including single-agent scenarios, multi-agent scenarios, and human-agent cooperation \cite{2309.07864}, \cite{2409.14457}. Effective communication capabilities, situational awareness, and the balance between autonomy and human control are identified as critical factors in developing collaborative agents that can work harmoniously with humans \cite{2401.03428}, \cite{2411.13904}. Moreover, the importance of evaluating these agents under realistic conditions is emphasized, with benchmarks such as TheAgentCompany providing a means to assess their performance on consequential real-world tasks \cite{2412.14161}.

However, the development and deployment of LLM-based agents also raise significant concerns regarding security vulnerabilities and privacy risks, particularly in multi-agent settings \cite{2308.05391}, \cite{2405.16310}. The potential for these agents to simulate human trust behavior and their implications on user trust dynamics in supply chains are areas of ongoing research \cite{2402.04559}, \cite{2405.16310}. Practical considerations for agentic LLM systems, including their design, deployment, and evaluation, are crucial for ensuring that these systems operate within acceptable boundaries of security, privacy, and ethics \cite{2412.04093}.

A deeper analysis of human-AI collaboration reveals common themes such as the emphasis on developing frameworks and models that facilitate effective interaction between humans and LLM-based agents \cite{2308.11432}, \cite{2409.14457}. The evaluation of these agents' performance, especially in scenarios involving cooperation and competition, is another prevalent theme, with studies employing various metrics and benchmarks to assess their capabilities \cite{2411.14033}, \cite{2308.05960}. Moreover, the exploration of LLM-based agents in diverse application domains, from travel planning to complex multi-agent systems, underscores their potential for wide-ranging impact \cite{2411.13904}, \cite{2402.16499}.

Despite the advancements, challenges such as scalability, complexity, and the need for explainability and transparency in LLM-based agents remain significant hurdles \cite{2403.11381}, \cite{2502.01390}. Addressing these challenges is essential for fostering trust and ensuring the safe and beneficial deployment of autonomous agents in real-world scenarios. The interplay between technical development, ethical considerations, and societal implications will continue to shape the trajectory of human-AI collaboration research, with LLMs at the forefront of this evolving landscape \cite{2412.15701}, \cite{2309.07864}. As research in this area continues to unfold, it is crucial to maintain a balanced perspective that acknowledges both the potential benefits and the challenges associated with the development and deployment of LLM-based collaborative agents.

\section{Evaluating LLMs and Agents: Benchmarks for Code Understanding, Generation, and Execution}
Evaluating LLMs and agents is crucial to understanding their capabilities in code understanding, generation, and execution. Recent studies have introduced various benchmarks to assess the performance of Large Language Models (LLMs) as agents in software engineering tasks. For instance, AgentBench \cite{2308.03688} provides a multi-dimensional benchmark to evaluate LLMs' reasoning and decision-making abilities in interactive environments, highlighting disparities in performance between commercial and open-source LLMs. Similarly, ML-Bench \cite{2311.09835} is designed to assess LLMs and agents on repository-level code understanding and generation tasks, revealing significant room for improvement in areas such as long code context interpretation and executable script generation.

The importance of training data quality and diversity has been emphasized in several studies, including \cite{2308.03688} and \cite{2311.09835}, which demonstrate the need for high-quality multi-turn alignment data to improve agent performance. In contrast, Agentless \cite{2407.01489} presents a simplistic approach to automated software development, achieving surprising performance on the SWE-bench Lite benchmark and suggesting that complex agent-based approaches may not be necessary. This highlights the ongoing debate about the complexity of agent-based approaches, with some studies advocating for simplicity \cite{2407.01489} and others introducing more complex benchmarks like AgentBench \cite{2308.03688} and ML-Bench \cite{2311.09835}.

The need for effective evaluation benchmarks is a common theme across these studies, with various methodologies employed to assess the capabilities of LLMs and agents. For example, DA-Code \cite{2410.07331} introduces an agent data science code generation benchmark for large language models, while LLMArena \cite{2402.16499} provides a dynamic multi-agent environment to assess the capabilities of large language models. These benchmarks have highlighted the challenges of interpreting long code contexts and generating executable scripts, as well as the importance of autonomy in software development.

The technical details and methodologies employed in these studies vary, with AgentBench \cite{2308.03688} consisting of 9,641 annotated examples across 18 GitHub repositories and ML-Bench \cite{2311.09835} employing two setups: ML-LLM-Bench for assessing LLMs' text-to-code conversion and ML-Agent-Bench for testing autonomous agents in end-to-end task execution. The evaluation metrics used also differ, with studies employing Pass@5 rates \cite{2308.03688}, success rates \cite{2311.09835}, and cost measurements \cite{2407.01489} to assess the performance of LLMs and agents.

Despite these developments, several limitations and challenges remain, including the hallucinated outputs and difficulties with bash script generation experienced by current LLMs \cite{2308.03688}. The complexity of agent-based approaches can also be a barrier to adoption, as highlighted by Agentless \cite{2407.01489}. Furthermore, evaluating the performance of LLMs and agents on repository-level code understanding and generation tasks is a challenging task, requiring large-scale benchmarks like ML-Bench \cite{2311.09835}. The need for high-quality training data and diverse evaluation scenarios is emphasized in several studies, including \cite{2308.03688} and \cite{2311.09835}, highlighting the ongoing challenges in developing effective LLM-based autonomous agents.

In conclusion, the evaluation of LLMs and agents is a critical aspect of understanding their capabilities in code understanding, generation, and execution. Recent studies have introduced various benchmarks and methodologies to assess the performance of LLMs as agents in software engineering tasks, highlighting the importance of training data quality, autonomy, and effective evaluation metrics. While limitations and challenges remain, these developments have significant implications for the future of software development and the role of LLM-based autonomous agents in this field, as discussed in \cite{2410.07331} and \cite{2402.16499}. As research continues to advance in this area, it is likely that we will see further improvements in the capabilities of LLMs and agents, leading to increased adoption and applications in software development, as envisioned in \cite{2308.03688} and \cite{2311.09835}.

\section{Technical Challenges and Open Problems in LLMs and Agent Development}
The development of Large Language Models (LLMs) has led to significant advancements in autonomous agent research, with various studies exploring the potential of LLM-based agents in achieving artificial general intelligence (AGI) \cite{2402.02716}, \cite{2308.11432}. However, several technical challenges and open problems remain to be addressed in this field. One of the primary concerns is the lack of autonomy and self-improvement in LLMs, which hinders their ability to interact with their environment and make decisions autonomously \cite{2409.14457}, \cite{2411.14033}. Furthermore, the planning and decision-making capabilities of LLM agents are still not well understood, with various approaches being proposed to improve their planning ability, including task decomposition, plan selection, and external modules \cite{2402.02716}, \cite{2407.01489}.

The integration of LLMs with autonomous agents has been identified as a key area of research, with the potential to achieve artificial collective intelligence through Large Language Model-based multi-agent systems (LaMAS) \cite{2411.14033}, \cite{2408.02479}. LaMAS offers several advantages, including dynamic task decomposition, higher flexibility, and proprietary data preservation, making it an attractive solution for various applications, including software engineering \cite{2409.02977}, \cite{2308.05960}. However, the development of LaMAS requires addressing technical challenges such as data privacy, business incentives, and evaluation metrics \cite{2412.04093}, \cite{2308.03688}.

The need for unified standards and benchmarking to qualify an LLM solution as an LLM-based agent has been emphasized in several studies \cite{2308.11432}, \cite{2407.01489}. The absence of standardized evaluation metrics and benchmarks hinders the comparison and assessment of different LLM-based agents, making it challenging to determine their effectiveness and potential applications \cite{2411.13768}, \cite{2409.14457}. Moreover, software engineering has been identified as a key application area for LLM-based agents, with tasks such as code generation, vulnerability detection, and test generation being explored \cite{2407.01489}, \cite{2308.05960}. However, the challenges associated with these applications, including the need for high-quality training data and the potential risks of biases and errors, must be addressed \cite{2409.02977}, \cite{2412.17481}.

In contrast to the development of complex autonomous software agents, a simplistic approach called Agentless has been proposed, which achieves high performance and low cost without using complex autonomous software agents \cite{2407.01489}. This approach challenges the need for advanced LLM-based agents and highlights the potential for alternative solutions in software engineering applications. However, other studies emphasize the importance of developing more advanced LLM-based agents that can interact with their environment and make decisions autonomously, underscoring the need for continued research in this area \cite{2402.02716}, \cite{2411.14033}.

The limitations and challenges associated with LLM-based agents must be addressed to realize their full potential. These include the lack of autonomy and self-improvement in LLMs, limited understanding of planning and decision-making in LLM agents, and the need for unified standards and benchmarking \cite{2409.14457}, \cite{2411.14033}. Furthermore, the challenges associated with software engineering applications, including code generation and vulnerability detection, must be addressed through continued research and development \cite{2407.01489}, \cite{2308.05960}. By exploring these challenges and limitations, researchers can develop more effective LLM-based agents that can be applied in various domains, ultimately contributing to the advancement of artificial general intelligence.

\section{Applications of LLMs and Agents: Daily Assistance, Code Development, and Social Simulations}
The applications of Large Language Models (LLMs) and LLM-based autonomous agents have been explored in various domains, including daily assistance, code development, and social simulations. According to \cite{2502.01390}, the use of LLM agents as daily assistants can be effective when employed in a plan-then-execute manner, but may also lead to mistrust if plans seem plausible but not well-executed. This highlights the importance of understanding user trust and team performance in human-agent collaboration, which is a recurring theme in the literature \cite{2408.02479, 2411.14033}. The potential benefits of LLMs and LLM-based agents in improving human-machine collaboration and decision-making are also emphasized in several papers \cite{2310.14985, 2402.01680}.

The application of LLM-based multi-agent systems has been proposed as a solution for achieving artificial collective intelligence, with advantages including dynamic task decomposition, higher flexibility, and proprietary data preservation \cite{2411.14033}. A survey of current practices and challenges in using LLMs and LLM-based agents for software engineering found that while LLMs have achieved remarkable success in areas like code generation and vulnerability detection, they also exhibit limitations and shortcomings that can be addressed by LLM-based agents \cite{2408.02479}. The survey on large language model based multi-agents provides an overview of the progress and challenges in this field, including the essential aspects of multi-agent systems based on LLMs and the mechanisms contributing to the growth of agents' capacities \cite{2402.01680}.

The technical aspects of LLMs and LLM-based agents are also explored in several papers, including the use of agent profiling and communication mechanisms in multi-agent systems \cite{2411.14033}. The importance of autonomy and self-improvement in LLM-based agents is highlighted in several papers, including the surveys on software engineering and multi-agents \cite{2408.02479, 2402.01680}. However, the papers also highlight several limitations and challenges associated with LLMs and LLM-based agents, including the lack of autonomy and self-improvement in LLMs, and the need for effective communication and coordination between agents \cite{2502.01390, 2411.14033}.

The implications of these developments are significant, as they suggest that LLMs and LLM-based agents have the potential to revolutionize various domains, including daily assistance, code development, and social simulations. However, they also highlight the need for further research into the challenges and limitations associated with these technologies \cite{2408.02479, 2411.14033}. The connections between these developments and other areas of research, such as human-computer interaction and artificial intelligence, are also worthy of exploration \cite{2310.14985, 2402.01680}. Overall, the applications of LLMs and LLM-based agents represent a rapidly evolving field, with significant potential for growth and development in the coming years \cite{2408.02479, 2411.14033}.

\section{Future Directions: Integrating LLMs with Multimodal Input, Edge AI, and Explainability}
As research in autonomous agents and multimodal large language models (MLLMs) continues to evolve, future directions point towards integrating LLMs with multimodal input, edge AI, and explainability to enhance their capabilities and applications. The concept of LLM-based Multi-Agent Systems (LaMAS), as introduced in \cite{2411.14033}, highlights the potential for achieving artificial collective intelligence through the integration of multiple agents and MLLMs. This development is further supported by surveys on large multimodal agents \cite{2402.15116} and autonomous agents through the lens of LLMs \cite{2404.04442}, which emphasize the importance of multimodal input in enhancing the capabilities of these systems.

The incorporation of edge AI into MLLMs is also a promising direction, as it enables real-time processing and decision-making in resource-constrained environments \cite{2310.02071}. This is particularly relevant for applications such as autonomous driving, where timely decisions are critical \cite{2311.12320}. Furthermore, the development of evaluation methodologies and benchmarks, such as PCA-EVAL \cite{2407.00118} and AgentBench \cite{2308.03688}, is essential for assessing the performance of autonomous agents and MLLMs.

Explainability is another crucial aspect of MLLMs, as it enables transparency and accountability in decision-making processes \cite{2402.11941}. Techniques such as prompting, reasoning, and tool utilization can be used to enhance the explainability of MLLMs \cite{2404.04442}. The introduction of frameworks like HOLMES \cite{2310.02071} and LaMAS \cite{2411.14033} provides a foundation for developing more complex and dynamic systems that leverage the capabilities of LLMs and MLLMs.

Despite these advancements, challenges remain, including the need for more advanced evaluation methodologies and benchmarks \cite{2408.15769}, as well as the difficulty in achieving artificial collective intelligence and abstract thinking in MLLMs \cite{2402.15527}. Additionally, the complexity of developing and integrating multiple AI systems, such as LaMAS and HOLMES, poses significant technical challenges \cite{2412.08442}. The potential limitations and biases of LLMs and MLLMs also need to be addressed to ensure their performance and decision-making capabilities are reliable and trustworthy \cite{2401.06805}.

In conclusion, the integration of LLMs with multimodal input, edge AI, and explainability is a promising direction for future research in autonomous agents and MLLMs. By addressing the challenges and limitations associated with these developments, researchers can create more advanced and dynamic systems that leverage the capabilities of LLMs and MLLMs to achieve artificial collective intelligence and enhance decision-making capabilities. The connections between these developments and their implications for various applications, including autonomous driving and multi-agent cooperation, highlight the significance of continued research in this area \cite{2401.06805, 2411.14033}. Ultimately, the advancement of autonomous agents and MLLMs will rely on the development of more sophisticated evaluation methodologies, frameworks, and techniques that enable transparency, accountability, and reliability in decision-making processes.

\section{Societal Implications and Ethical Considerations of Widespread LLM and Agent Adoption}
The widespread adoption of Large Language Models (LLMs) and autonomous agents has significant societal implications and raises important ethical considerations. As highlighted in \cite{2411.14033}, LLM-based multi-agent systems have the potential to perceive, control, and interact with their environment, which can be beneficial in various applications, but also poses risks if not properly regulated. A comprehensive survey on LLM-based autonomous agents \cite{2308.11432} provides an overview of the field, discussing construction, applications, and evaluation strategies, while also emphasizing the need for safety measures and regulations to mitigate potential risks.

The risks associated with scientific LLM agents are a major concern, as identified in \cite{2402.04247}, which highlights vulnerabilities to misuse and the need for safeguarding over autonomy. Furthermore, the ontological characterization of LLMs is investigated in \cite{2407.10735}, concluding that they do not meet necessary and sufficient conditions for autonomous agency, which has significant implications for their development and deployment. In contrast, other papers, such as \cite{2310.14985}, demonstrate the potential benefits of LLM-based agents in navigating dynamic social interactions, highlighting the importance of considering the societal implications of their adoption.

The need for safety measures and regulations to mitigate risks associated with LLM-based autonomous agents is a recurring theme across the literature \cite{2402.04247, 2407.10735}. The potential benefits of LLM-based autonomous agents in various applications, including scientific research and social interactions, are also highlighted \cite{2411.14033, 2308.11432, 2310.14985}. However, the concept of agency and autonomy in LLMs is a common thread, with some papers arguing that LLMs do not meet the necessary conditions for autonomous agency \cite{2407.10735}, while others assume that they can be considered autonomous agents or have the potential to become so \cite{2411.14033, 2308.11432}.

The paper \cite{2502.02649} presents a strong argument against the development of fully autonomous AI agents, highlighting the potential risks and consequences of such systems. In contrast, other papers, such as \cite{2409.14457}, discuss the state-of-the-art in large model-based agents, including cooperation paradigms, and highlight the potential benefits of these systems. The importance of considering the ethical implications of LLM-based autonomous agents is also emphasized in \cite{2402.04247}, which highlights the need for prioritizing safeguarding over autonomy in scientific research.

The technical details and methodologies employed in the development of LLM-based autonomous agents are also an important consideration. Papers such as \cite{2310.14985} employ frameworks for understanding social behaviors, while others, such as \cite{2411.14033}, provide an overview of the technical aspects of LLM-based multi-agent systems. However, the scalability and complexity of these systems are identified as potential limitations \cite{2411.14033}, highlighting the need for further research and development.

In conclusion, the societal implications and ethical considerations of LLM-based autonomous agents are complex and multifaceted. While these systems have the potential to bring significant benefits in various applications, they also pose risks if not properly regulated. As highlighted in \cite{2407.10735}, the concept of agency and autonomy in LLMs is still not well understood, highlighting the need for further philosophical and technical investigation. Ultimately, the development and deployment of LLM-based autonomous agents must be carefully considered, with a focus on prioritizing safety, ethics, and responsibility.

\section{Conclusion: The State of the Art in Large Language Models and Agents, and Paths Forward}
The state of the art in large language models and agents is characterized by significant advancements and ongoing debates about the nature of autonomous agency. As highlighted in \cite{2407.10735}, current large language models do not meet the conditions for autonomous agency due to their lack of individuality, normativity, and interactional symmetry, instead being viewed as interlocutors or linguistic automatons. In contrast, papers such as \cite{2308.11432} and \cite{2309.07870} introduce frameworks and libraries that support the development of autonomous language agents with capabilities such as planning, memory, tool usage, and multi-agent communication, underscoring the potential of large language models to contribute significantly to the development of autonomous agents.

A key theme emerging from the literature is the importance of interdisciplinary approaches, combining insights from computer science, philosophy, and cognitive sciences \cite{2407.16190}. The creation of accessible frameworks and tools, like the Agents library \cite{2309.07870}, is seen as crucial for advancing the field by making it easier for researchers and developers to build and experiment with autonomous language agents. Furthermore, theoretical models, such as those presented in \cite{2407.16190}, provide a threshold conception for artificial agents, suggesting that certain architectural elements and modules could pave the way for realizing agency in an artificial manner.

Despite these advancements, challenges persist in achieving autonomous agency, including issues related to embodiment, normativity, and true autonomy \cite{2308.11432}. Practical considerations, such as handling stochasticity and managing resources efficiently, also pose significant challenges when designing agentic large language models for real-world applications \cite{2412.04093}. The philosophical and ethical implications of developing autonomous agents based on large language models are becoming increasingly pressing, with questions about their development and deployment requiring careful consideration \cite{2407.16190}.

The literature also reveals contrasting viewpoints on what constitutes an autonomous agent and how achievable this goal is with current large language models. For instance, \cite{2407.10735} offers a philosophical critique of large language models' capacity for agency, while papers like \cite{2412.04093} focus on the practical aspects of implementing agentic large language model systems. Additionally, some researchers are optimistic about the potential of current architectures and modules to be adapted or combined to achieve artificial agency \cite{2403.11381}, while others are more skeptical about the current state of large language models in meeting the criteria for true autonomous agency \cite{2407.10735}.

Technical details and methodologies vary across the papers, with proposals for different architectural elements and modules that could enhance the autonomy of large language model-based agents \cite{2412.04093}. The introduction of libraries like Agents indicates a move towards more accessible and standardized tools for developing autonomous language agents \cite{2309.07870}. Theoretical modeling, as employed in \cite{2407.16190}, provides a conceptual framework for artificial agency, suggesting a pathway for future research and development.

In conclusion, the research into autonomous agents based on large language models is characterized by significant progress, ongoing debates, and persistent challenges. As the field continues to evolve, it is essential to address the philosophical, ethical, and practical implications of developing autonomous agents, ensuring that these systems are aligned with human values and promote beneficial outcomes \cite{2407.16190}. By acknowledging the complexities and limitations of current large language models, researchers can work towards creating more advanced and truly autonomous systems that have the potential to transform various aspects of society \cite{2308.11432}. Ultimately, the development of autonomous agents based on large language models requires a comprehensive and interdisciplinary approach, incorporating insights from computer science, philosophy, and cognitive sciences to create systems that are not only technologically sophisticated but also ethically responsible and socially beneficial \cite{2407.16190}.

\section{References and Further Reading}
For further exploration of autonomous agents in Multi-Layer Large Language Models (MLLM), several key references provide foundational insights and cutting-edge research directions. The concept of LLM-based Multi-Agent Systems, as discussed in \cite{2411.14033}, offers a promising avenue for achieving artificial collective intelligence by leveraging the strengths of both large language models and multi-agent systems. This is complemented by comprehensive surveys such as \cite{2308.11432}, which outline frameworks, applications, and challenges associated with LLM-based autonomous agents, emphasizing their potential in diverse fields and the need for unified evaluation strategies.

Delving deeper into the operational aspects of these agents, research like \cite{2402.02716} categorizes planning in LLM agents, discussing task decomposition, plan selection, and memory integration, thereby shedding light on the complexities of agent decision-making processes. Furthermore, efforts to formally specify the behavior of LLM-based agents, as proposed in \cite{2310.08535}, aim to simplify their design and implementation by allowing for high-level behavioral specifications. This technical approach is crucial for advancing the autonomy and adaptability of these agents.

The ontological status of large language models as autonomous agents is a topic of debate, with analyses like \cite{2407.10735} questioning whether LLMs can be considered truly autonomous under embodied theories of mind, instead positing them as sophisticated linguistic tools capable of structured but non-purposeful engagement. This philosophical inquiry underscores the importance of defining autonomy and agency in the context of artificial intelligence, a theme that resonates throughout the literature on LLM-based agents.

The development of benchmarks and frameworks for assessing and orchestrating LLM-augmented autonomous agents, such as those described in \cite{2308.05960} and \cite{2309.07870}, highlights the community's recognition of the need for standardized tools and methodologies to propel this field forward. Moreover, surveys and analyses like \cite{2409.14457} and \cite{2404.13501} provide overviews of the state-of-the-art in large model-based agents, including their cooperative capabilities, security concerns, and memory mechanisms, which are essential for understanding both the potential and the limitations of these systems.

The capability of LLM-augmented autonomous agents to cooperate is a critical aspect of their functionality, with studies like \cite{2403.11381} evaluating their performance in multi-agent environments. Additionally, theoretical works such as \cite{2407.16190} and practical considerations outlined in \cite{2412.04093} contribute to a nuanced understanding of artificial agency in the context of large language models, emphasizing both the advancements achieved and the challenges that remain.

For readers interested in the broader implications of autonomous agents in MLLM, including their potential applications in software engineering as discussed in \cite{2404.13501}, and their ethical considerations, delving into the referenced literature will provide a comprehensive insight into the current research landscape. The interplay between technical innovation, philosophical inquiry, and practical application is at the heart of advancing our understanding of autonomous agents in MLLM, with works like \cite{2308.11432}, \cite{2402.02716}, and \cite{2310.08535} serving as foundational references for future research directions.

In conclusion, the study of autonomous agents in Multi-Layer Large Language Models is a vibrant and multidisciplinary field, encompassing technical, philosophical, and practical dimensions. Through the exploration of the referenced works, including \cite{2411.14033}, \cite{2308.11432}, \cite{2402.02716}, \cite{2310.08535}, \cite{2407.10735}, \cite{2308.05960}, \cite{2309.07870}, \cite{2409.14457}, \cite{2404.13501}, \cite{2403.11381}, \cite{2407.16190}, and \cite{2412.04093}, researchers and practitioners can gain a deeper understanding of the current state of knowledge, the challenges that lie ahead, and the potential pathways for innovation in this exciting area of artificial intelligence research.

\section{Conclusion}
The concept of autonomous agents in Multi-Large Language Models (MLLM) has garnered significant attention in recent years, with a plethora of research papers exploring its potential applications and challenges. Based on our analysis of relevant papers, including \cite{2411.14033}, \cite{2308.11432}, \cite{2402.02716}, \cite{2407.16190}, \cite{2407.10735}, \cite{2407.01489}, \cite{2404.13501}, \cite{2012.12600}, \cite{2310.08535}, \cite{2405.06643}, \cite{2308.05960}, \cite{2409.14457}, \cite{2402.01680}, \cite{2309.07870}, \cite{2405.02957}, \cite{2405.11106}, \cite{2403.11381}, \cite{2412.04093}, \cite{2412.03359}, and \cite{2407.07086}, we summarize the key findings, common themes, contrasting viewpoints, technical details, and limitations of this emerging field.

The development of LLM-based autonomous agents has shown remarkable potential in achieving human-level intelligence, with applications in various fields such as social science, natural science, and engineering, as highlighted in \cite{2411.14033} and \cite{2308.11432}. The integration of LLMs with other modules or tools to enhance their autonomous capabilities is a recurring theme across the papers, with \cite{2402.02716} and \cite{2407.16190} emphasizing the importance of environment-interaction properties, dynamic task decomposition, and organic specialization in achieving autonomy. Furthermore, the potential of LaMAS to provide higher flexibility, proprietary data preserving, and feasibility of monetization for each entity is highlighted in \cite{2409.14457} and \cite{2308.05960}.

However, there are also contrasting viewpoints on the concept of autonomy in MLLM, with some papers arguing that LLMs are not yet autonomous agents, but rather interlocutors or linguistic automatons, lacking individuality, normativity, and interactional asymmetry, as discussed in \cite{2407.10735} and \cite{2407.01489}. Others propose that LLMs can be considered as agents with a dynamic framework of factors influencing their actions and goals, as suggested in \cite{2405.06643} and \cite{2310.08535}. The role of embodiment in achieving autonomy is also debated, with some papers emphasizing the importance of sensorimotor and biological embodiment, while others argue that textual embodiment and computational embodiment can suffice, as seen in \cite{2403.11381} and \cite{2412.04093}.

From a technical perspective, the use of external tools and modules to enhance LLM capabilities is a common approach, with \cite{2404.13501} and \cite{2012.12600} exploring the development of LaMAS protocols considering technical requirements, data privacy, and business incentives. Evaluation strategies for LLM-based autonomous agents are also being explored, including metrics for assessing planning ability and decision-making, as discussed in \cite{2308.05960} and \cite{2409.14457}. Nevertheless, the lack of a clear definition of autonomy in the context of MLLM is identified as a significant challenge, with \cite{2402.02716} and \cite{2407.16190} emphasizing the need for more research on the obstacles to building artificial agents with LLMs.

In conclusion, our analysis highlights the significance of autonomous agents in MLLM, with key themes and developments including the integration of LLMs with other modules, the importance of environment-interaction properties, and the potential of LaMAS. However, there are also limitations and challenges to be addressed, such as the lack of a clear definition of autonomy and the need for more research on evaluation strategies and potential risks. Future research should focus on developing more robust and flexible LaMAS protocols, exploring new evaluation strategies, and investigating the potential implications and connections of this emerging field, ultimately contributing to the advancement of autonomous agents in MLLM.

\end{document}